% Rôles des membres et usage de l'IA : partie commune aux deux rapports (% Rôles des membres et usage de l'IA : partie commune aux deux rapports (% Rôles des membres et usage de l'IA : partie commune aux deux rapports (% Rôles des membres et usage de l'IA : partie commune aux deux rapports (\input{roles_et_ia}).

\section*{Contributions des membres}

\begin{table}[H]\centering\small
\begin{tabular}{p{3.2cm}p{12cm}}\toprule
Membre & Ce qu'il a fait \\\midrule
\textbf{Yann Linganzi} &
\textbf{Conception et conduite du projet.} Choix de la question (acheter, rester neutre ou vendre le Nikkei, couvert
ou non contre le yen) et de sa traduction en deux classifications ; collecte des données (versions 1 et 2), choix des
sources et des décalages de publication ; audit des fuites ; exploration et construction des variables ; découpage
chronologique ; pipeline, trois modèles et réglage Optuna ; choix de la log-loss comme critère ; calibration de Platt
et règle de décision par espérance de gain, figées avant le test ; évaluation unique du test ; explicabilité
(tranches, permutation, SHAP, dépendances partielles) ; annexe mathématique (logistique codée à la main, sensibilité,
moindres carrés) ; seconde version 1999-2026 avec le contrôle des séries Bloomberg ; explorations après coup et
protocole figé de la troisième version ; API, tableau de bord et site ; rédaction des rapports. \\[3pt]
\textbf{arnshl} &
Captures du terminal Bloomberg (huit écrans : inflation japonaise, fondamentaux et décomposition du rendement du
Nikkei, révisions des bénéfices, rendement des bénéfices contre le taux américain à 10 ans, trajectoires du bénéfice
par action), avec les vérifications visuelles et les pistes de variables qu'elles suggèrent ; chiffres qui justifient
les conclusions des salles III et VII du site ; relecture et fusion des contributions sur GitHub. \\[3pt]
\textbf{jovi-mce} &
Version anglaise du site (\texttt{site/en/}) et sa mise en ligne. \\\bottomrule
\end{tabular}\end{table}

\section*{Usage de l'IA}

La réflexion et la méthode sont les nôtres. Yann Linganzi a posé le problème et décidé de chaque choix de méthode :
les variables et leurs dates de publication, le découpage dans le temps, les modèles comparés, le critère, la
calibration, la règle de décision et l'engagement pris avant d'ouvrir le test, puis, après coup, les vérifications à
mener et le protocole de la troisième version. Il a testé chacun de ces choix et validé chaque résultat : audit des
fuites, vérifications qui retrouvent à l'identique la décision et l'évaluation figées, tests automatiques, lecture et
interprétation des chiffres.

Un assistant IA (Claude, d'Anthropic) a servi d'aide à la programmation : écriture du code Python des modules et des
notebooks à partir de ces choix, mise en forme du site et des figures, et aide à la relecture et à la mise en forme de
la rédaction. Tout ce qu'il a produit a été relu, exécuté et vérifié avant d'être gardé ; quand un résultat ne
tenait pas, c'est la méthode qui a été corrigée, pas le chiffre. La soutenance est assurée par le groupe.
).

\section*{Contributions des membres}

\begin{table}[H]\centering\small
\begin{tabular}{p{3.2cm}p{12cm}}\toprule
Membre & Ce qu'il a fait \\\midrule
\textbf{Yann Linganzi} &
\textbf{Conception et conduite du projet.} Choix de la question (acheter, rester neutre ou vendre le Nikkei, couvert
ou non contre le yen) et de sa traduction en deux classifications ; collecte des données (versions 1 et 2), choix des
sources et des décalages de publication ; audit des fuites ; exploration et construction des variables ; découpage
chronologique ; pipeline, trois modèles et réglage Optuna ; choix de la log-loss comme critère ; calibration de Platt
et règle de décision par espérance de gain, figées avant le test ; évaluation unique du test ; explicabilité
(tranches, permutation, SHAP, dépendances partielles) ; annexe mathématique (logistique codée à la main, sensibilité,
moindres carrés) ; seconde version 1999-2026 avec le contrôle des séries Bloomberg ; explorations après coup et
protocole figé de la troisième version ; API, tableau de bord et site ; rédaction des rapports. \\[3pt]
\textbf{arnshl} &
Captures du terminal Bloomberg (huit écrans : inflation japonaise, fondamentaux et décomposition du rendement du
Nikkei, révisions des bénéfices, rendement des bénéfices contre le taux américain à 10 ans, trajectoires du bénéfice
par action), avec les vérifications visuelles et les pistes de variables qu'elles suggèrent ; chiffres qui justifient
les conclusions des salles III et VII du site ; relecture et fusion des contributions sur GitHub. \\[3pt]
\textbf{jovi-mce} &
Version anglaise du site (\texttt{site/en/}) et sa mise en ligne. \\\bottomrule
\end{tabular}\end{table}

\section*{Usage de l'IA}

La réflexion et la méthode sont les nôtres. Yann Linganzi a posé le problème et décidé de chaque choix de méthode :
les variables et leurs dates de publication, le découpage dans le temps, les modèles comparés, le critère, la
calibration, la règle de décision et l'engagement pris avant d'ouvrir le test, puis, après coup, les vérifications à
mener et le protocole de la troisième version. Il a testé chacun de ces choix et validé chaque résultat : audit des
fuites, vérifications qui retrouvent à l'identique la décision et l'évaluation figées, tests automatiques, lecture et
interprétation des chiffres.

Un assistant IA (Claude, d'Anthropic) a servi d'aide à la programmation : écriture du code Python des modules et des
notebooks à partir de ces choix, mise en forme du site et des figures, et aide à la relecture et à la mise en forme de
la rédaction. Tout ce qu'il a produit a été relu, exécuté et vérifié avant d'être gardé ; quand un résultat ne
tenait pas, c'est la méthode qui a été corrigée, pas le chiffre. La soutenance est assurée par le groupe.
).

\section*{Contributions des membres}

\begin{table}[H]\centering\small
\begin{tabular}{p{3.2cm}p{12cm}}\toprule
Membre & Ce qu'il a fait \\\midrule
\textbf{Yann Linganzi} &
\textbf{Conception et conduite du projet.} Choix de la question (acheter, rester neutre ou vendre le Nikkei, couvert
ou non contre le yen) et de sa traduction en deux classifications ; collecte des données (versions 1 et 2), choix des
sources et des décalages de publication ; audit des fuites ; exploration et construction des variables ; découpage
chronologique ; pipeline, trois modèles et réglage Optuna ; choix de la log-loss comme critère ; calibration de Platt
et règle de décision par espérance de gain, figées avant le test ; évaluation unique du test ; explicabilité
(tranches, permutation, SHAP, dépendances partielles) ; annexe mathématique (logistique codée à la main, sensibilité,
moindres carrés) ; seconde version 1999-2026 avec le contrôle des séries Bloomberg ; explorations après coup et
protocole figé de la troisième version ; API, tableau de bord et site ; rédaction des rapports. \\[3pt]
\textbf{arnshl} &
Captures du terminal Bloomberg (huit écrans : inflation japonaise, fondamentaux et décomposition du rendement du
Nikkei, révisions des bénéfices, rendement des bénéfices contre le taux américain à 10 ans, trajectoires du bénéfice
par action), avec les vérifications visuelles et les pistes de variables qu'elles suggèrent ; chiffres qui justifient
les conclusions des salles III et VII du site ; relecture et fusion des contributions sur GitHub. \\[3pt]
\textbf{jovi-mce} &
Version anglaise du site (\texttt{site/en/}) et sa mise en ligne. \\\bottomrule
\end{tabular}\end{table}

\section*{Usage de l'IA}

La réflexion et la méthode sont les nôtres. Yann Linganzi a posé le problème et décidé de chaque choix de méthode :
les variables et leurs dates de publication, le découpage dans le temps, les modèles comparés, le critère, la
calibration, la règle de décision et l'engagement pris avant d'ouvrir le test, puis, après coup, les vérifications à
mener et le protocole de la troisième version. Il a testé chacun de ces choix et validé chaque résultat : audit des
fuites, vérifications qui retrouvent à l'identique la décision et l'évaluation figées, tests automatiques, lecture et
interprétation des chiffres.

Un assistant IA (Claude, d'Anthropic) a servi d'aide à la programmation : écriture du code Python des modules et des
notebooks à partir de ces choix, mise en forme du site et des figures, et aide à la relecture et à la mise en forme de
la rédaction. Tout ce qu'il a produit a été relu, exécuté et vérifié avant d'être gardé ; quand un résultat ne
tenait pas, c'est la méthode qui a été corrigée, pas le chiffre. La soutenance est assurée par le groupe.
).

\section*{Contributions des membres}

\begin{table}[H]\centering\small
\begin{tabular}{p{3.2cm}p{12cm}}\toprule
Membre & Ce qu'il a fait \\\midrule
\textbf{Yann Linganzi} &
\textbf{Conception et conduite du projet.} Choix de la question (acheter, rester neutre ou vendre le Nikkei, couvert
ou non contre le yen) et de sa traduction en deux classifications ; collecte des données (versions 1 et 2), choix des
sources et des décalages de publication ; audit des fuites ; exploration et construction des variables ; découpage
chronologique ; pipeline, trois modèles et réglage Optuna ; choix de la log-loss comme critère ; calibration de Platt
et règle de décision par espérance de gain, figées avant le test ; évaluation unique du test ; explicabilité
(tranches, permutation, SHAP, dépendances partielles) ; annexe mathématique (logistique codée à la main, sensibilité,
moindres carrés) ; seconde version 1999-2026 avec le contrôle des séries Bloomberg ; explorations après coup et
protocole figé de la troisième version ; API, tableau de bord et site ; rédaction des rapports. \\[3pt]
\textbf{arnshl} &
Captures du terminal Bloomberg (huit écrans : inflation japonaise, fondamentaux et décomposition du rendement du
Nikkei, révisions des bénéfices, rendement des bénéfices contre le taux américain à 10 ans, trajectoires du bénéfice
par action), avec les vérifications visuelles et les pistes de variables qu'elles suggèrent ; chiffres qui justifient
les conclusions des salles III et VII du site ; relecture et fusion des contributions sur GitHub. \\[3pt]
\textbf{jovi-mce} &
Version anglaise du site (\texttt{site/en/}) et sa mise en ligne. \\\bottomrule
\end{tabular}\end{table}

\section*{Usage de l'IA}

La réflexion et la méthode sont les nôtres. Yann Linganzi a posé le problème et décidé de chaque choix de méthode :
les variables et leurs dates de publication, le découpage dans le temps, les modèles comparés, le critère, la
calibration, la règle de décision et l'engagement pris avant d'ouvrir le test, puis, après coup, les vérifications à
mener et le protocole de la troisième version. Il a testé chacun de ces choix et validé chaque résultat : audit des
fuites, vérifications qui retrouvent à l'identique la décision et l'évaluation figées, tests automatiques, lecture et
interprétation des chiffres.

Un assistant IA (Claude, d'Anthropic) a servi d'aide à la programmation : écriture du code Python des modules et des
notebooks à partir de ces choix, mise en forme du site et des figures, et aide à la relecture et à la mise en forme de
la rédaction. Tout ce qu'il a produit a été relu, exécuté et vérifié avant d'être gardé ; quand un résultat ne
tenait pas, c'est la méthode qui a été corrigée, pas le chiffre. La soutenance est assurée par le groupe.
